% Generated draft: requires mathematical and source review.
\documentclass[11pt]{article}
\usepackage[margin=25mm]{geometry}
\usepackage{fontspec}
\setmainfont{Latin Modern Roman}
\usepackage{amsmath,amssymb,mathtools,unicode-math}
\setmathfont{Latin Modern Math}
\usepackage{enumitem,needspace,xurl,hyperref}
\setlength{\parindent}{0pt}
\setlength{\parskip}{5pt}
\setlength{\emergencystretch}{3em}
\newcommand{\N}{\mathbb{N}}
\newcommand{\Z}{\mathbb{Z}}
\newcommand{\Q}{\mathbb{Q}}
\newcommand{\R}{\mathbb{R}}
\newcommand{\C}{\mathbb{C}}
\newcommand{\F}{\mathbb{F}}
\newcommand{\A}{\mathbb{A}}
\newcommand{\PP}{\mathbb{P}}
\newcommand{\E}{\mathbb{E}}
\newcommand{\eps}{\varepsilon}
\newcommand{\dd}{\,\mathrm{d}}
\begin{document}
\section*{UnsolvedMath ID: 700983}
\textbf{Source problem code:} catalogue entry
\subsection{Definitions and mathematical statement}
For a finite group \(G\), define
\[
\pi(G)=\{p\in\N:p\text{ is prime and }p\mid |G|\}.
\]
A group is simple if it is nontrivial and has no normal subgroups other than itself and the identity subgroup. For a positive integer \(n\), a finite simple group \(G\) is called a \(K_n\)-group if \(|\pi(G)|=n\).

Let
\[
\mathcal K_4=\{[G]_{\cong}:G\text{ is a finite simple group and }|\pi(G)|=4\},
\]
where \([G]_{\cong}\) denotes the isomorphism class of \(G\). The question is whether \(\mathcal K_4\) is finite or infinite. Counting up to isomorphism is the conventional interpretation adopted here; the supplied statement does not explicitly specify this convention.

The supplied background states that there are eight \(K_3\)-groups, that the \(K_4\)-groups have been classified subject to the classification of finite simple groups (CFSG), and that further significant results appear in the cited 2001 paper. These are source-reported background assertions, not independently verified claims in this draft.
\subsection{Short English statement}
Are there finitely or infinitely many isomorphism classes of finite simple groups whose orders have exactly four distinct prime divisors? The supplied registry labels this problem open; that label is not an independently verified assessment.
\subsection{Sources}
\begin{itemize}
\item Local UnsolvedMath registry, \texttt{data/unsolvedmath/problems.json} and \texttt{catalog/problems.json}; supplied complete record with permanent ID \(700983\), source code \texttt{catalogue entry}, title \emph{Kourovka Notebook Problem 13.65}.
\item W. J. Shi, proposer, 1995. The supplied attribution identifies Problem 13.65 of \emph{The Kourovka Notebook (Unsolved Problems in Group Theory)}, 13th issue (1995), as reproduced in the 21st issue edition (2026), arXiv:1401.0300v48. The supplied collection description names E. I. Khukhro and V. D. Mazurov as editors.
\item W. J. Shi, contribution in \emph{Group Theory in China}, \emph{Mathematics and Its Applications}, vol. 365, Kluwer, 1996, pp. 163--181. Supplied as the reference for classification of \(K_4\)-groups subject to CFSG; no contribution title is supplied.
\item Yann Bugeaud, Zhenfu Cao, and M. Mignotte, \emph{Journal of Algebra}, vol. 241 (2001), pp. 658--668. Supplied as a reference for significant further results; no article title is supplied.
\end{itemize}
\end{document}
